\section{No equivariant map from the two-sphere}\label{sec:coindex}

The lower bounds for the chromatic number in~\cite{Z} and in Section~\ref{sec:complete} show that the box complex of a graph admits no equivariant map \emph{to} a sphere of low dimension. For graphs with a quantum $3$-coloring, maps in the other direction are restricted as well. The \emph{coindex} of a space $X$ with a free involution is the largest $k$ for which there is an equivariant map from the sphere $S^k$ with the antipodal involution to $X$~\cite{Mat}.

\begin{proposition}\label{prop:coindex}
Let $G$ be a graph with a quantum $3$-coloring. Then there is no equivariant map from the sphere $S^2$ with the antipodal involution to $\B(G)$. If $G$ is not bipartite, the coindex of $\B(G)$ is one.
\end{proposition}

The proof imports a theorem of Simons, Tardif and Wehlau, which we state first. Let $H$ be a graph. An \emph{arc} of $H$ is an ordered pair $(a,b)$ of adjacent vertices. Let $\Gamma(H)$ be the group generated by the arcs of $H$ with the relations
\begin{equation}\label{eq:stw-relation}
 (a,b)(c,b)^{-1}(c,d)(a,d)^{-1}=1\qquad\text{for all vertices with }a\sim b\sim c\sim d\sim a .
\end{equation}
If $a=c$ or $b=d$ the relation holds in the free group, so it does not matter whether such degenerate cycles are included. For a closed walk $v_0,v_1,\ldots,v_n=v_0$ of odd length $n$ in $H$, its \emph{signature} is the element
\[
 \prod_{i=0}^{n-1}(v_{2i},v_{2i+1})\,(v_{2i+2},v_{2i+1})^{-1}
\]
of $\Gamma(H)$, with the indices read modulo $n$ and the factors in the order of increasing $i$. In~\cite{STW} the box complex is realized as the order complex of the poset of pairs $(A,B)$ of nonempty sets of vertices with every vertex of $A$ adjacent to every vertex of $B$, with the involution $(A,B)\mapsto(B,A)$. The cells of $\B(H)$ in Section~\ref{sec:box} are indexed by this poset, so the order complex is the barycentric subdivision of $\B(H)$, and the two spaces are equivariantly homeomorphic.

\begin{theorem}[Simons, Tardif and Wehlau {\cite[Theorem~1 with Lemma~4, arXiv version~2]{STW}}]\label{thm:stw}
If there is an equivariant map from $S^2$ to $\B(H)$, then $H$ has a closed walk of odd length whose signature is the identity of $\Gamma(H)$.
\end{theorem}

\begin{proof}[Proof of Proposition~\ref{prop:coindex}]
Index the colors by $\mathbb Z/3$ and let $(P_{u,a})$ be a quantum $3$-coloring of $G$ on $\C^d$. By~(Q1) the projectors at adjacent vertices commute. For an arc $(u,v)$ and $\theta\in\R$ put
\[
 U_\theta(u,v)=\sum_{a\in\mathbb Z/3}\bigl(e^{i\theta}P_{u,a}P_{v,a+1}+e^{-i\theta}P_{u,a}P_{v,a-1}\bigr).
\]
The nine products $P_{u,a}P_{v,b}$ are pairwise orthogonal projections with sum $I$, and the three with $a=b$ vanish by the edge rule. So $U_\theta(u,v)$ is unitary, and $U_\theta(v,u)=U_\theta(u,v)^{-1}$, because $P_{u,a}P_{v,a+1}$ has the phase $e^{i\theta}$ in $U_\theta(u,v)$ and $e^{-i\theta}$ in $U_\theta(v,u)$.

\emph{The product along two edges depends only on the ends.} Let $u\sim v\sim w$, where $u=w$ is allowed. We claim that
\begin{equation}\label{eq:two-edges}
 U_\theta(u,v)\,U_\theta(v,w)=\sum_{a\in\mathbb Z/3}P_{u,a}\bigl(P_{w,a}+e^{-2i\theta}P_{w,a+1}+e^{2i\theta}P_{w,a-1}\bigr).
\end{equation}
Indeed, the left side is the sum of the terms $\varphi(b-a)\varphi(c-b)\,P_{u,a}P_{v,b}P_{w,c}$ over all $a,b,c$ with $b\ne a$ and $c\ne b$, where $\varphi(1)=e^{i\theta}$ and $\varphi(-1)=e^{-i\theta}$. If $c=a$, the two values $b=a\pm1$ occur, both with phase $1$, and $P_{u,a}(P_{v,a+1}+P_{v,a-1})P_{w,a}=P_{u,a}(I-P_{v,a})P_{w,a}=P_{u,a}P_{w,a}$, because $P_{u,a}P_{v,a}=0$. If $c=a+1$, only $b=a-1$ occurs, with phase $e^{-2i\theta}$, and $P_{u,a}P_{v,a-1}P_{w,a+1}=P_{u,a}(I-P_{v,a}-P_{v,a+1})P_{w,a+1}=P_{u,a}P_{w,a+1}$, because $P_{u,a}P_{v,a}=0$ and $P_{v,a+1}P_{w,a+1}=0$. The case $c=a-1$ is symmetric, with phase $e^{2i\theta}$.

\emph{A representation of $\Gamma(G)$.} By~\eqref{eq:two-edges}, for $a\sim b\sim c\sim d\sim a$ the products $U_\theta(a,b)U_\theta(b,c)$ and $U_\theta(a,d)U_\theta(d,c)$ are equal. As $U_\theta(c,b)^{-1}=U_\theta(b,c)$ and $U_\theta(a,d)^{-1}=U_\theta(d,a)$, the assignment $(u,v)\mapsto U_\theta(u,v)$ respects the relations~\eqref{eq:stw-relation} and defines a homomorphism from $\Gamma(G)$ to the unitary group of $\C^d$.

\emph{Conclusion.} Suppose that there is an equivariant map $S^2\to\B(G)$. By Theorem~\ref{thm:stw} there is a closed walk $v_0,\ldots,v_n=v_0$ of odd length with trivial signature. The homomorphism sends the factor $(v_{2i},v_{2i+1})(v_{2i+2},v_{2i+1})^{-1}$ of the signature to $U_\theta(v_{2i},v_{2i+1})U_\theta(v_{2i+1},v_{2i+2})$. As $i$ runs from $0$ to $n-1$ and the indices are read modulo the odd number $n$, the walk is traversed twice. So the signature is sent to $H(\theta)^2$, where
\[
 H(\theta)=U_\theta(v_0,v_1)\,U_\theta(v_1,v_2)\cdots U_\theta(v_{n-1},v_n),
\]
and $H(\theta)^2=I$ for every $\theta$. Thus $H(\theta)$ is a unitary involution, $(I-H(\theta))/2$ is a projection that depends continuously on $\theta$, and its rank is constant. For $\theta=2\pi/3$ put $Z_u=\sum_ae^{2\pi ia/3}P_{u,a}$. Then $U_\theta(u,v)=Z_u^{-1}Z_v$ on every arc, so the product telescopes and $H(2\pi/3)=Z_{v_0}^{-1}Z_{v_n}=I$. For $\theta=\pi$ every $U_\theta(u,v)$ equals $-I$, so $H(\pi)=(-1)^nI=-I$. The rank of $(I-H(\theta))/2$ is $0$ in the first case and $d$ in the second, a contradiction.

Finally, if $G$ contains an odd cycle $C$, then $\B(C)$ is a circle on which the exchange of the shores is the half-turn~\cite[proof of Theorem~7.1]{Z}, and the inclusion $\B(C)\subseteq\B(G)$ is an equivariant map from $S^1$. So the coindex is at least one, and it is at most one by the first part.
\end{proof}

The representation $U_\theta$ interpolates between two homomorphisms that every quantum $3$-coloring provides, the telescoping one at $\theta=2\pi/3$ and the sign of the length at $\theta=\pi$. The proposition restricts the graphs with a quantum $3$-coloring from the side opposite to the chromatic number: their box complexes can have large height, as the graphs of~\cite{Z} show, but their coindex is at most one.
